% Part: many-valued-logic
% Chapter: three-valued-logics
% Section: goedel

\documentclass[../../../include/open-logic-section]{subfiles}

\begin{document}

\olfileid{mvl}{thr}{god}

\olsection{منطقات غودل}

وضع كورت غودل سلسلة من المنطقات ذات~$\OLId{latin-ordinary}{n}$ من القيم؛ يشتمل كل واحد
منها على جميع~!!{formula}s الصحيحة حدسيًا، ويشتمل عليه المنطق
الكلاسيكي. وهذا أول أفراد السلسلة غير التافهة:

\begin{defn}\ollabel{defn:goedel}
مصفوفة \emph{منطق غودل الثلاثي القيم}~$\LogGod$ هي:
\begin{enumerate}
  \item لغة القضايا القياسية~$\Lang L_\OLMathNumeral{0}$ ذات الروابط
  $\lfalse$ و$\lnot$ و$\land$ و$\lor$ و$\lif$.
  \item مجموعة قيم الصدق~$\OLId{latin-looped}{V} = \{\True, \Undef, \False\}$.
  \item القيمة المميزة وحدها هي~$\True$؛ أي~$\OLId{latin-looped}{V}^+ = \{\True\}$.
  \item في الثابت~$\lfalse$ يكون~$\tf{\lfalse} = \False$؛ وأما
  دوال صدق سائر الروابط فتعطيها الجداول:
  \begin{center}
    \begin{tabular}{c|c}
      $\tf{\lnot}[\LogGod]$ & \\
      \hline
      $\True$ & $\False$ \\
      $\Undef$ & $\False$ \\
      $\False$ & $\True$
    \end{tabular}
    \quad
    \begin{tabular}{c|ccc}
      $\tf{\land}[\LogGod]$ & $\True$ & $\Undef$ & $\False$ \\
      \hline
      $\True$ & $\True$ & $\Undef$ & $\False$ \\
      $\Undef$ & $\Undef$ & $\Undef$ & $\False$\\
      $\False$ & $\False$ & $\False$ & $\False$
    \end{tabular}
    \\[2ex]
    \begin{tabular}{c|ccc}
      $\tf{\lor}[\LogGod]$ & $\True$ & $\Undef$ & $\False$ \\
      \hline
      $\True$ & $\True$ & $\True$ & $\True$ \\
      $\Undef$ & $\True$ & $\Undef$ & $\Undef$ \\
      $\False$ & $\True$ & $\Undef$ & $\False$
    \end{tabular}
    \quad
    \begin{tabular}{c|ccc}
      $\tf{\lif}[\LogGod]$ & $\True$ & $\Undef$ & $\False$ \\
      \hline
      $\True$ & $\True$ & $\Undef$ & $\False$ \\
      $\Undef$ & $\True$ & $\True$ & $\False$  \\
      $\False$ & $\True$ & $\True$ & $\True$
    \end{tabular}
  \end{center}
\end{enumerate}
\end{defn}

وجدولا~$\land$ و$\lor$ هنا هما الجدولان نفسيهما في منطق~\L
ukasiewicz وفي منطق كليني القوي؛ وإنما يختلف~$\lnot$ و$\lif$ في
هذه المنطقات. ففي منطق غودل~$\tf{\lnot}(\Undef) = \False$.
وكذلك، بخلاف منطق~\L ukasiewicz ومنطق كليني، يكون
$\tf{\lif}(\Undef, \False) = \False$؛ وأما
$\tf{\lif}(\Undef, \Undef) = \True$، فيخالف منطق كليني ويوافق
منطق~\L ukasiewicz.

وارتباط هذا المنطق بالمنطق الحدسي ظاهر في~$\LogGod[\OLMathNumeral{3}]$: فكل حقيقة
حدسية تحصيل في~$\LogGod[\OLMathNumeral{3}]$، وكثير من التحصيلات الكلاسيكية التي لا تصح حدسيًا لا
تكون تحصيلات في~$\LogGod[\OLMathNumeral{3}]$ أيضًا. فمن أمثلتها هذه الصيغ:
\begin{align*}
  & \OLId{latin-ordinary}{p} \lor \lnot \OLId{latin-ordinary}{p} && (\OLId{latin-ordinary}{p} \lif \OLId{latin-ordinary}{q}) \lif (\lnot \OLId{latin-ordinary}{p} \lor \OLId{latin-ordinary}{q}) \\
  & \lnot\lnot \OLId{latin-ordinary}{p} \lif \OLId{latin-ordinary}{p} && \lnot(\lnot \OLId{latin-ordinary}{p} \land \lnot \OLId{latin-ordinary}{q}) \lif (\OLId{latin-ordinary}{p} \lor \OLId{latin-ordinary}{q}) \\
  & ((\OLId{latin-ordinary}{p} \lif \OLId{latin-ordinary}{q}) \lif \OLId{latin-ordinary}{p}) \lif \OLId{latin-ordinary}{p} && \lnot(\OLId{latin-ordinary}{p} \lif \OLId{latin-ordinary}{q}) \lif (\OLId{latin-ordinary}{p} \land \lnot \OLId{latin-ordinary}{q})
\end{align*}
ومع ذلك توجد تحصيلات لـ~$\LogGod[\OLMathNumeral{3}]$ لا تصح حدسيًا، ومنها
$\lnot\lnot \OLId{latin-ordinary}{p} \lor \lnot \OLId{latin-ordinary}{p}$ و~$(\OLId{latin-ordinary}{p} \lif \OLId{latin-ordinary}{q}) \lor (\OLId{latin-ordinary}{q} \lif \OLId{latin-ordinary}{p})$.

\begin{prob}
  ضع جداول صدق تثبت أن الصيغ الآتية تحصيلات في~$\LogGod[\OLMathNumeral{3}]$:
  \begin{align*}
    & \lnot\lnot \OLId{latin-ordinary}{p} \lor \lnot \OLId{latin-ordinary}{p}\\
    & (\OLId{latin-ordinary}{p} \lif \OLId{latin-ordinary}{q}) \lor (\OLId{latin-ordinary}{q} \lif \OLId{latin-ordinary}{p}) \\
    & \lnot(\OLId{latin-ordinary}{p} \land \OLId{latin-ordinary}{q}) \lif (\lnot \OLId{latin-ordinary}{p} \lor \lnot \OLId{latin-ordinary}{q}) \\
    & (\OLId{latin-ordinary}{p} \lif \OLId{latin-ordinary}{q}) \lor (\OLId{latin-ordinary}{q} \lif \OLId{latin-ordinary}{r}) \lor (\OLId{latin-ordinary}{r} \lif \OLId{latin-ordinary}{s})
  \end{align*}
\end{prob}

\begin{prob}
  ضع جداول صدق تثبت أن الصيغ الآتية ليست تحصيلات في~$\LogGod[\OLMathNumeral{3}]$:
  % تصحيح أسلوبي (0396-N1): استُعمل فعل الأمر بدل المصدر.
  \begin{align*}
    & (\OLId{latin-ordinary}{p} \lif \OLId{latin-ordinary}{q}) \lif (\lnot \OLId{latin-ordinary}{p} \lor \OLId{latin-ordinary}{q}) \\
    & \lnot(\lnot \OLId{latin-ordinary}{p} \land \lnot \OLId{latin-ordinary}{q}) \lif (\OLId{latin-ordinary}{p} \lor \OLId{latin-ordinary}{q}) \\
    & ((\OLId{latin-ordinary}{p} \lif \OLId{latin-ordinary}{q}) \lif \OLId{latin-ordinary}{p}) \lif \OLId{latin-ordinary}{p} \\
    & \lnot(\OLId{latin-ordinary}{p} \lif \OLId{latin-ordinary}{q}) \lif (\OLId{latin-ordinary}{p} \land \lnot \OLId{latin-ordinary}{q})
  \end{align*}
\end{prob}

\begin{prob}
  أي العلاقات الآتية تصح في منطق غودل؟ ضع جدول صدق لكل واحدة.
  \begin{enumerate}
    \item $\OLId{latin-ordinary}{p}, \OLId{latin-ordinary}{p} \lif \OLId{latin-ordinary}{q} \Entails \OLId{latin-ordinary}{q}$
    \item $\OLId{latin-ordinary}{p} \lor \OLId{latin-ordinary}{q}, \lnot \OLId{latin-ordinary}{p} \Entails \OLId{latin-ordinary}{q}$
    \item $\OLId{latin-ordinary}{p} \land \OLId{latin-ordinary}{q} \Entails \OLId{latin-ordinary}{p}$
    \item $\OLId{latin-ordinary}{p} \Entails \OLId{latin-ordinary}{p} \land \OLId{latin-ordinary}{p}$
    \item $\OLId{latin-ordinary}{p} \Entails \OLId{latin-ordinary}{p} \lor \OLId{latin-ordinary}{q}$
  \end{enumerate}
\end{prob}

\end{document}
